\section{Citologia}

Schema generale dei tipi di cellula: cellula batterica, cellula vegetale, cellula animale. Oltre alle cellule esistono anche batteriofagi, virus, viroidi e prioni.

% ---------------------------------------------------------------------------
\subsection{Virus e batteriofagi}

\begin{itemize}
  \item dimensioni: 10--300 nm;
  \item \textbf{virus}: si riproducono all'interno di cellule eucariotiche;
  \item \textbf{batteriofagi} (o fagi): si riproducono all'interno dei batteri;
  \item utilizzano l'apparato sintetico della cellula ospite per sintetizzare e assemblare le molecole che li compongono.
\end{itemize}

\rubrica{Struttura del virus}
\begin{itemize}
  \item \textbf{genoma virale}: DNA o RNA;
  \item \textbf{capside}: involucro proteico costituito da subunità (capsomeri) disposte in modo regolare; protegge il genoma virale in ambiente extracellulare e consente la penetrazione del virus nella cellula ospite;
  \item \textbf{pericapside} (o peplos): ulteriore involucro lipoproteico, talvolta presente all'esterno del capside.
\end{itemize}

\rubrica{Classificazione dei virus}
\begin{itemize}
  \item in base al tipo di cellula infettata: batteriofagi/fagi, virus vegetali, virus animali;
  \item in base all'acido nucleico: desossiribovirus (virus a DNA), ribovirus o retrovirus (virus a RNA);
  \item in base alla forma: allungati, sferici, poliedrici;
  \item in base al tipo di simmetria: cubica, elicoidale, complessa.
\end{itemize}

% ---------------------------------------------------------------------------
\subsection{Replicazione di virus e batteriofagi}

\rubrica{Ciclo litico dei batteriofagi}
\begin{enumerate}
  \item aggancio: il fago aderisce alla superficie cellulare del batterio;
  \item penetrazione: il DNA del fago entra nella cellula batterica;
  \item replicazione e sintesi: il DNA del fago viene replicato e le sue proteine sintetizzate;
  \item assemblaggio: le componenti del fago vengono assemblate in nuovi virus;
  \item rilascio: la cellula batterica si lisa liberando molti fagi, che possono infettare altre cellule.
\end{enumerate}

\rubrica{Ciclo lisogenico dei batteriofagi}
\begin{enumerate}
  \item aggancio;
  \item penetrazione;
  \item integrazione: il DNA fagico si integra nel DNA batterico (\textit{profago});
  \item replicazione: il profago integrato si replica quando viene replicato il DNA batterico; le cellule figlie possono esibire nuove proprietà.
\end{enumerate}
Il fago \textbf{lambda} è un esempio di fago temperato, che può effettuare il ciclo litico o quello lisogenico.

\rubrica{Replicazione a confronto}
Nel \textbf{virus}: adesione a parete o membrana cellulare $\rightarrow$ rilascio di enzimi che consentono l'adesione alla cellula ospite $\rightarrow$ dissoluzione del capside $\rightarrow$ ingresso del genoma virale $\rightarrow$ duplicazione del genoma $\rightarrow$ sintesi delle proteine del capside $\rightarrow$ assemblaggio delle nuove particelle virali.

Nel \textbf{batteriofago}: adesione alla parete batterica $\rightarrow$ iniezione dell'acido nucleico nel batterio $\rightarrow$ duplicazione del genoma $\rightarrow$ sintesi delle proteine del capside $\rightarrow$ assemblaggio dei nuovi fagi $\rightarrow$ lisi della cellula ospite $\rightarrow$ fuoriuscita dei fagi, che possono infettare altri batteri.

% ---------------------------------------------------------------------------
\subsection{Viroidi e prioni}

\rubrica{Viroidi}
\begin{itemize}
  \item piccole molecole di RNA chiuse ad anello, capaci di autoreplicazione;
  \item un solo acido nucleico;
  \item privi di capside;
  \item infettano preferenzialmente le cellule vegetali.
\end{itemize}

\rubrica{Prioni}
\begin{itemize}
  \item molecole proteiche normalmente presenti in soggetti sani;
  \item in particolari condizioni alterano la propria configurazione spaziale;
  \item trasmettono la propria configurazione alterata alle molecole sane.
\end{itemize}

% ---------------------------------------------------------------------------
\subsection{Cellula procariote}

Componenti: parete cellulare, membrana plasmatica, capsula, DNA (nel nucleoide/area nucleare), plasmidi, ribosomi, citoplasma, mesosoma, pili, flagello, granuli di riserva.

\rubrica{Classificazione dei batteri per forma}
\begin{itemize}
  \item \textbf{bacilli}: a bastoncino;
  \item \textbf{cocchi}: a sfera;
  \item \textbf{diplococchi}: due cocchi uniti;
  \item \textbf{streptococchi}: disposti a catena;
  \item \textbf{stafilococchi}: disposti a grappolo;
  \item \textbf{spirilli}: a spirale;
  \item \textbf{vibrioni}: a virgola;
  \item \textbf{spirochete}: con più curve.
\end{itemize}

\rubrica{Altre classificazioni dei batteri}
\begin{itemize}
  \item in base alla temperatura: \textit{criofili/psicrofili} (\textit{Pseudomonas}, \textit{Lactobacillus}), \textit{mesofili} (\textit{Salmonella}, stafilococco, streptococco), \textit{termofili} (\textit{Thermus aquaticus});
  \item in base alla presenza/assenza di ossigeno: aerobi, anaerobi, aerobi/anaerobi facoltativi;
  \item in base al grado di acidità: acidofili, alcalofili;
  \item in base alla concentrazione salina: alofili (elevate concentrazioni saline).
\end{itemize}

\rubrica{Parete cellulare batterica}
Costituita da uno strato di peptidoglicano.
\begin{itemize}
  \item \textbf{Gram-positivi}: parete spessa (spesso strato di peptidoglicano); la capsula può essere presente;
  \item \textbf{Gram-negativi}: parete sottile, con una membrana esterna oltre allo strato di peptidoglicano.
\end{itemize}

\rubrica{Duplicazione dei procarioti}
Avviene per scissione binaria: il cromosoma (DNA), collegato al mesosoma, si duplica a partire da un'origine di replicazione (replicazione bidirezionale); le due origini migrano verso i poli della cellula; la membrana si accresce verso l'interno e si forma una nuova parete $\rightarrow$ due cellule figlie che si separano.

% ---------------------------------------------------------------------------
\subsection{Cellula eucariote}

Organuli: nucleo (involucro nucleare, nucleoplasma, nucleolo, cromatina), membrana plasmatica, reticolo endoplasmatico (rugoso e liscio), ribosomi, complesso di Golgi, lisosomi, perossisomi, mitocondri, citoscheletro (microfilamenti, microtubuli, filamenti intermedi), centrioli, vescicole, citosol. La cellula \textbf{vegetale} presenta in più parete cellulare, cloroplasti, vacuolo centrale e plasmodesmi.

\rubrica{Membrana plasmatica}
Doppio strato fosfolipidico (fosfolipidi con testa idrofilica e coda idrofobica); costituisce il limite esterno del citoplasma; presenta canali proteici di trasporto per le sostanze idrosolubili, che non possono attraversare la regione fosfolipidica.

\rubrica{Involucro nucleare}
Sistema di due membrane concentriche attraversate da pori nucleari, che facilitano il trasporto di molecole tra nucleo e citoplasma; sulla superficie dell'involucro (membrana esterna, affacciata al citoplasma) sono presenti ribosomi.

\rubrica{Reticolo endoplasmatico}
\begin{itemize}
  \item \textbf{RE rugoso}: con ribosomi sulla superficie (sede della sintesi proteica);
  \item \textbf{RE liscio}: privo di ribosomi.
\end{itemize}

\rubrica{Complesso di Golgi}
Insieme di cisterne; riceve vescicole originate dal RE (modificazione e smistamento delle proteine) ed emette vescicole in fase di gemmazione.

\rubrica{Mitocondri}
Membrana esterna e membrana interna (che si ripiega formando le creste); matrice interna e compartimento intermembrana; sede del metabolismo energetico.

\rubrica{Cloroplasti}
Membrana esterna e membrana interna; tilacoidi impilati in \textit{grana} (granum); stroma (fluido interno).

\rubrica{Componenti cellulari a confronto}
Presenza (X) dei principali componenti nei procarioti (Eubatteri, Archea) e negli eucarioti (Protisti, Funghi, Piante, Animali).
\begin{table}[htbp]
  \centering \footnotesize \renewcommand{\arraystretch}{1.15}
  \begin{tabular}{|l|c|c|c|c|c|c|}
    \hline
    \textbf{Componente} & \textbf{Eubatt.} & \textbf{Archea} & \textbf{Protisti} & \textbf{Funghi} & \textbf{Piante} & \textbf{Animali} \\
    \hline
    Nucleoide            & X & X &        &        &   &   \\ \hline
    Nucleo               &   &   & X      & X      & X & X \\ \hline
    Involucro nucleare   &   &   & X      & X      & X & X \\ \hline
    Nucleolo             &   &   & X      & X      & X & X \\ \hline
    Membrana plasmatica  & X & X & X      & X      & X & X \\ \hline
    Parete cellulare     & X & X & alcuni & X      & X &   \\ \hline
    Ribosomi             & X & X & X      & X      & X & X \\ \hline
    Reticolo endoplasm.  &   &   & X      & X      & X & X \\ \hline
    Complesso di Golgi   &   &   & X      & X      & X & X \\ \hline
    Lisosoma             &   &   & X      & alcuni & X & X \\ \hline
    Mitocondrio          &   &   & X      & X      & X & X \\ \hline
    Cloroplasto          &   &   & alcuni &        & X &   \\ \hline
    Vacuolo centrale     &   &   &        &        & X &   \\ \hline
    Microfilamenti       &   &   & X      & X      &   & X \\ \hline
    Microtubuli          &   &   & X      & X      & X & X \\ \hline
    Filamenti intermedi  &   &   & X      &        &   & X \\ \hline
  \end{tabular}
\end{table}
